%%%%%%%%%%%%%%%%%%%%%%%%%%%%%%%%%%%%%%%%%
% Их сургуулийн оюутны тезис  
% LaTeX Загвар
% Version 2.3 (25/3/16)
%
% Энэ загвар нь дараах сайтаас авсан загварын монгол хувилбар юм.
% http://www.LaTeXTemplates.com
%
% Version 2.x major modifications by:
% Vel (vel@latextemplates.com)
%
% Анхдагч загварын эх үүсвэр:
% Steve Gunn (http://users.ecs.soton.ac.uk/srg/softwaretools/document/templates/)
% Sunil Patel (http://www.sunilpatel.co.uk/thesis-template/)
%
% Загварын лиценз:
% CC BY-NC-SA 3.0 (http://creativecommons.org/licenses/by-nc-sa/3.0/)
%
%%%%%%%%%%%%%%%%%%%%%%%%%%%%%%%%%%%%%%%%%

%-------------------------------------------------------------------------------
%	PACKAGES AND OTHER DOCUMENT CONFIGURATIONS
%-------------------------------------------------------------------------------

\documentclass[
12pt, % Баримтын фонтын хэмжээ, сонголт: 10pt, 11pt, 12pt
oneside, % Хоёр талаар хэвлэж үдэхээр тохируулсан. Нэг тал бол комментыг арилга
%chapterinoneline,% Нэг мөрөнд бүлгийн дугаар, нэрийг гаргах
english, % babel багцын хэлний тохиргоо
onehalfspacing, % Мөр хоорондын зай. Сонголтууд: singlespacing, onehalfspacing, doublespacing
%draft, % Ноорог горимд шилжихийн тулд комментыг арилга(зураг, холбоос, hboxes гарахгүй)
nolistspacing, % Хэрэв мөр хоорондын зай onehalfspacing эсвэл doublespacing бол, жагсаалтын мөр хоорондын зайг single болгохын тулд комментыг арилга
%liststotoc, % Зураг/хүснэгт/бусад жагсаалтыг гарчигт оруулахын тулд комментыг арилга
%toctotoc, % Uncomment to add the main table of contents to the table of contents
%parskip, % Параграф хооронд зай оруулахын тулд комментыг арилга
%nohyperref, % hyperref багцыг ачаалахгүй бол комментыг арилга
headsepline, % Толгой мөрийн доогуур шугам татахын тулд комментыг арилга
]{MUST-Thesis} % Энэ класс файл нь баримтын бүтцийг тодорхойлно

\usepackage[utf8]{inputenc} % Олон улсын тэмдэгт оруулахад хэрэгтэй
\usepackage[T2A]{fontenc} % Олон улсын тэмдэгтийн гаралтын кодчилол
\usepackage[mongolian]{babel}

\usepackage{titlesec}
\usepackage[table]{xcolor}
\usepackage{rotating} % эргүүлэх
\usepackage{ragged2e}
\let\oldtabular\tabular
\let\endoldtabular\endtabular
\renewenvironment{tabular}
   {\bgroup\singlespacing\oldtabular}%
   {\endoldtabular\egroup}
\let\oldverse\verse
\let\endoldverse\endverse
\renewenvironment{verse}
   {\bgroup\singlespacing\oldverse}%
   {\endoldverse\egroup}

\usepackage{caption,longtable,ltcaption}
\DeclareCaptionJustification{nohyphen}{\hyphenpenalty=10000}
\captionsetup{justification=nohyphen}

\usepackage{pgf}
\usepackage{tikz} % зураг зурах
\usetikzlibrary{shapes,arrows,automata}

\definecolor{OliveGreen}{cmyk}{0.64,0,0.95,0.40}
\definecolor{CadetBlue}{cmyk}{0.62,0.57,0.23,0}
\definecolor{Cover}{RGB}{255,255,255}
% \definecolor{Cover}{RGB}{102,255,255}
% \definecolor{Chapter}{RGB}{255,220,180}

\definecolor{lightlightgray}{gray}{0.9}
% \usepackage{listings}
\usepackage{listings,lstautogobble}
\renewcommand{\lstlistingname}{Эх код} % Програмын эх код хэвлэх
\lstset{
    language=C,                             % Code langugage
    basicstyle=\linespread{1.0}\ttfamily,                   % Code font
    keywordstyle=\color{OliveGreen},        % Keywords font ('*' = uppercase)
    commentstyle=\color{CadetBlue},         % Comments font
    numbers=left,                           % Line nums position
    numberstyle=\tiny,                      % Line-numbers fonts
    stepnumber=1,                           % Step between two line-numbers
    numbersep=5pt,                          % How far are line-numbers from code
    backgroundcolor=\color{lightlightgray}, % Choose background color
    frame=none,                             % A frame around the code
    tabsize=2,                              % Default tab size
    captionpos=t,                           % Caption-position = bottom
    breaklines=true,                        % Automatic line breaking?
    breakatwhitespace=false,                % Automatic breaks only at whitespace?
    showspaces=false,                       % Dont make spaces visible
    showtabs=false,                         % Dont make tabls visible
    columns=flexible,                       % Column format
    morekeywords={__global__, __device__},  % CUDA specific keywords
}

% source code
\colorlet{punct}{red!60!black}
\definecolor{background}{HTML}{EEEEEE}
\definecolor{delim}{RGB}{20,105,176}
\colorlet{numb}{magenta!60!black}
\lstdefinelanguage{json}{
    basicstyle=\normalfont\ttfamily,
    numbers=left,
    numberstyle=\scriptsize,
    stepnumber=1,
    numbersep=8pt,
    showstringspaces=false,
    breaklines=true,
    frame=lines,
    autogobble=true,
    backgroundcolor=\color{background},
    literate=
     *{0}{{{\color{numb}0}}}{1}
      {1}{{{\color{numb}1}}}{1}
      {2}{{{\color{numb}2}}}{1}
      {3}{{{\color{numb}3}}}{1}
      {4}{{{\color{numb}4}}}{1}
      {5}{{{\color{numb}5}}}{1}
      {6}{{{\color{numb}6}}}{1}
      {7}{{{\color{numb}7}}}{1}
      {8}{{{\color{numb}8}}}{1}
      {9}{{{\color{numb}9}}}{1}
      {:}{{{\color{punct}{:}}}}{1}
      {,}{{{\color{punct}{,}}}}{1}
      {\{}{{{\color{delim}{\{}}}}{1}
      {\}}{{{\color{delim}{\}}}}}{1}
      {[}{{{\color{delim}{[}}}}{1}
      {]}{{{\color{delim}{]}}}}{1},
}

\usepackage[autostyle=false]{csquotes} % Ном зүйд хэлнээс хамаарсан хашилт оруулахад хэрэгтэй

\usepackage[backend=bibtex,natbib=true,sorting=none,sortcites]{biblatex} % Ном зүйд bibtex -г ашиглах

\usepackage{float}
\usepackage{geometry}
\usepackage{graphicx}
\usepackage{subcaption}
\usepackage{hyperref}
\usepackage{indentfirst}
\usepackage{enumitem}
\usepackage{textcase}
\usepackage{tocloft}
\usepackage{multicol}
\usepackage{pdflscape}
\usepackage{etoolbox}
\usepackage{amsmath}
\usepackage{multirow}
\usepackage{setspace}
\graphicspath{ {./} }


\newcommand{\authorshipname}{Зохиогч эрхийн хамгаалал}
\newcommand{\abbrevname}{Товчилсон үгс}
\newcommand{\constantsname}{Ашигласан томьёонууд}
\newcommand{\symbolsname}{Таних тэмдэгтүүд}
\newcommand{\acknowledgementname}{Талархал}
\newcommand{\abstractname}{Хураангуй}
 
%-------------------------------------------------------------------------------
%	THESIS INFORMATION
%-------------------------------------------------------------------------------

\thesistitle{Шаблом засах} % Таны ажлын нэр, нүүр болон хураангуй хуудсанд ашигласан. Өөр газарт бол \ttitle командыг хэрэглэнэ
\thesistype{Бакалаврын төслийн ажил} % Удирдагчийн нэр, нүүр хуудсанд ашиглана. Дурын газарт бол \thesisname командыг хэрэглэнэ
\supervisor{Магистр Т.Золбоо} % Удирдагчийн нэр, нүүр хуудсанд ашиглана. Дурын газарт бол \supname командыг хэрэглэнэ
%\reader{Доктор (PhD) Ч.Цэнд-Аюуш} % Шүүмжлэгчийн нэр, Дурын газарт бол \readname командыг хэрэглэнэ
%\advisor{Магистр Б.Гүндсамбуу, Магистр Г.Цэнд-Аюуш} % Зөвлөгчийн нэр, Дурын газарт бол \advicename командыг хэрэглэнэ
\degreeind{D0613030000000014} % Мэргэжлийн индекс, нүүр болон хураангуй хуудсанд ашигласан. Өөр газарт бол \degreeid командыг хэрэглэнэ
\degree{Мэдээллийн систем} % Боловсролын зэрэг, нүүр болон хураангуй хуудсанд ашигласан. Өөр газарт бол \degreename командыг хэрэглэнэ
\authorshort{Д.Мөнхтүвшин} % Таны товч нэр, нүүр болон хураангуй хуудсанд ашигласан. Дурын газарт бол \shortname командыг хэрэглэнэ 
\authorlong{Давхарбаярын Мөнхтүвшин} % Таны бүтэн нэр, нүүр болон хураангуй хуудсанд ашигласан. Дурын газарт бол \longname командыг хэрэглэнэ 
\addresses{B150920009@mymust.net} % Таны хаяг, одоогоор ашиглаагүй. Өөр газарт бол \addressname командыг хэрэглэнэ
\subject{Мэдээллийн технологи} % Таны салбар, одоогоор ашиглаагүй. Дурынр газарт бол \subjectname командыг ашиглана
\keywords{шаблом} % Түлхүүр үгс, одоогоор ашиглаагүй. Дурын газарт бол \keywordnames командыг хэрэглэнэ
\university{\href{http://www.must.edu.mn}{Шинжлэх Ухаан Технологийн Их Сургууль}} % Их сургуулийн нэр ба веб хаяг. Дурын газарт бол \univname командыг хэрэглэнэ 
\department{Мэдээллийн технологийн салбар} % Сургууль/тэнхмийн нэр, нүүр болон хураангуй хуудсанд ашигласан. Дурын газарт бол\deptname командыг хэрэглэнэ
\deptchair{Проф. PhD. А.Эрдэнэбаатар} % Тэнхим/эрхлэгчийн нэр, нүүр болон хураангуй хуудсанд ашигласан. Дурын газарт бол\chairname командыг хэрэглэнэ
%\group{Робот техникийн баг} % Судалгааны баг/тэнхмийн нэр, нүүр хуудсанд ашигласан. Дурын газарт бол \groupname командыг хэрэглэнэ
\faculty{\href{http://www.sict.edu.mn}{Мэдээлэл, Холбооны Технологийн Сургууль}} % Салбар сургууль/факультетийн нэр, нүүр болон хураангуй хуудсанд ашигласан. Дурын газарт бол \facname командыг ашиглана

\hypersetup{pdftitle=\ttitle} % Pdf файлын гарчиг
\hypersetup{pdfauthor=\shortname} % Pdf файлын зохиогчийн нэр
\hypersetup{pdfkeywords=\keywordnames} % Pdf файлын түлхүүр үгс
\hypersetup{allcolors=black} % Pdf файлын бүх холбоос хар өнгөтэй

\renewcommand{\thesection}{\thechapter.\number\numexpr\value{section}-1\relax}
\renewcommand{\thesubsection}{\thesection.\number\numexpr\value{subsection}-1\relax}
\renewcommand{\thesubsubsection}{\thesubsection.\number\numexpr\value{subsubsection}-1\relax}
\setcounter{secnumdepth}{3}

\setcounter{chapter}{1}

\begin{document}

 % Агуулгын өмнөх хуудас дугаарлалт: i, ii, iii, iv... г.м.

\pagestyle{plain} % Тезисийн загварыг дуудах хүртэлх толгой мөрийн суурь загвар
\frontmatter
%-------------------------------------------------------------------------------
%	TITLE PAGE
%-------------------------------------------------------------------------------
\pagecolor{Cover}
\begin{titlepage}
\begin{center}

{\scshape\LARGE \univname\par} % Их сургуулийн нэр
{\scshape\Large \facname\par}\vspace{0.5cm} % Их сургуулийн нэр

\begin{figure}[!htbp]
\centering
\includegraphics[scale=0.2]{Figures/MUST_logo.png}
\end{figure}

\vspace{1cm}
\hfill \large{\longname} \\

\vspace{1cm}

{\huge \bfseries \ttitle\par}\vspace{0.4cm} % Тезисийн нэр

\vspace{3cm}
\textsc{\Large {\thesisname}}\\ % Тезисийн төрөл

\vfill

\large {Улаанбаатар хот} \\
 
\end{center}
\end{titlepage}
\pagecolor{white}

%-------------------------------------------------------------------------------
%	SUBTITLE PAGE
%-------------------------------------------------------------------------------

\begin{titlepage}
\begin{center}

{\scshape\LARGE \univname\par} % Их сургуулийн нэр
{\scshape\Large \facname\par}\vspace{0.5cm} % Их сургуулийн нэр

\vspace{2cm}
\hfill \large{\deptname} \\

\vspace{2cm}

{\huge \bfseries \ttitle\par}\vspace{0.4cm} % Тезисийн нэр

\vspace{2cm}

\begin{minipage}[t] {0.9\textwidth}
\begin{flushleft} 
\normalsize

Мэргэжлийн индекс: \degreeid \\
Мэргэжил: \degreename \\[2cm]

\emph{Удирдагч:} {\supname} \\% Удирдагчийн нэр
%\emph{Зөвлөгч:} {\advicename} \\ % Зөвлөгч нарын нэрс
\emph{Гүйцэтгэгч:} {\shortname} \\ % Зохиогчийн нэр

\end{flushleft}
\end{minipage}

\vfill

\large {Улаанбаатар хот} \\
{\large 2018 он}\\ % Date

\end{center}
\end{titlepage}

  % Нүүр хуудас
%-------------------------------------------------------------------------------
%	ABSTRACT PAGE
%-------------------------------------------------------------------------------

\addchaptertocentry{Хураангуй} % Хураангуйг гарчигт нэмэх

\begin{center}
{\scshape\Large \univname\par} % Их сургуулийн нэр
{\scshape\large \facname\par}\vspace{0.5cm} % Их сургуулийн нэр
{\huge\textbf{{Хураангуй}} \par}
\bigskip
{\Large{\ttitle} \par} % Тезисийн нэр
\bigskip

{\normalsize \shortname \par} % Зохиогчийн нэр
\addressname
\end{center}

\textit{\textbf{Түлхүүр үгс: \keywordnames}}
\bigskip

Төслийн ажлын тухай хураангуйг энд бичих



 % Ажлын хураангуй
%-------------------------------------------------------------------------------
%	ACKNOWLEDGEMENTS
%-------------------------------------------------------------------------------

\begin{acknowledgements}
\addchaptertocentry{\acknowledgementname}

Юуны түрүүнд энэхүү төслийг гүйцэлдүүлэх явцыг удирдан явуулсан Шинжлэх Ухаан Технологийн Их Сургуулийн харьяа Мэдээлэл, Холбооны Технологийн Сургуулийн Магистр Ж.Золжаргал-д гүнээ талархал илэрхийлж байна. Мөн тус сургуулийн Компьютерийн ухааны салбарын Профессор (PhD) А.Эрдэнэбаатар, Доктор (PhD), дэд профессор Б.Батзолбоо, Доктор (PhD) А.Хүдэр, Доктор (PhD) Ч.Цэнд-Аюуш, Доктор (PhD) Ч.Золзаяа, Доктор (PhD) Д.Энхжаргал, Доктор (PhD) Ч.Эрдэнэбат, Магистр Т.Золбоо, Магистр Д.Золбоо, Магистр Б.Сод-Од, Магистр Б.Гүндсамбуу, Магистр Б.\\Мөнхбуян, Магистр Г.Цэнд-Аюуш, Магистр Б.Намжилдорж тэргүүтэй багш, профессоруудад талархал илэрхийлж байна.

\end{acknowledgements}

 % Талархлын хуудас
%-------------------------------------------------------------------------------
%	ABBREVIATIONS
%-------------------------------------------------------------------------------

\begin{abbreviations}{ll} % Товчлолын жагсаалт оруулах (хоёр багатай хүснэгт)
\addchaptertocentry{\abbrevname}

\textbf{AI} & \textbf{A}rtifical \textbf{I}ntelligence\\
\textbf{ХОУ} & \textbf{Х}иймэл \textbf{О}юун \textbf{У}хаан\\
\textbf{АНУ} & \textbf{А}мерикийн \textbf{Н}эгдсэн \textbf{У}лс\\
\textbf{JSON} & \textbf{J}ava\textbf{S}cript \textbf{O}bject \textbf{N}otation\\
\textbf{API} & \textbf{A}pplication \textbf{P}rogramming \textbf{I}nterface\\
\textbf{CPU} & \textbf{C}entral \textbf{P}rocessing \textbf{U}nit\\
\textbf{RAM} & \textbf{R}andom \textbf{A}ccess \textbf{M}emory\\


\end{abbreviations}

 % Товчилсон үгс
\tableofcontents % Гарчиг хэвлэх
\newpage
\listoffigures % Зургийн жагсаалт хэвлэх
\newpage
\listoftables % Хүснэгтийн жагсаалт хэвлэх

\mainmatter


\section{Удиртгал}

    \subsection{Зорилго}

    \subsection{Систем хөгжүүлэх үндэслэл}
    
    \subsection{Судлагдсан байдал}

    \subsection{Шинэлэг тал}

    \subsection{Системийн ач холбогдол}

    \subsection{Системийн хамрах хүрээ}
% Бүлэг 1


%-------------------------------------------------------------------------------

% Агуулгад ашигласан хэвшүүлэлтийн зарим командын тодорхойлолт
\newcommand{\keyword}[1]{\textbf{#1}}
\newcommand{\tabhead}[1]{\textbf{#1}}
\newcommand{\code}[1]{\texttt{#1}}
\newcommand{\file}[1]{\texttt{\bfseries#1}}
\newcommand{\option}[1]{\texttt{\itshape#1}}

\section{Судалгаа}	
\subsection{Системийн судалгаа}
	Энэ төрлийн ситемүүд нь computer-ийн ухааны computer vision-ы салбар дээр үндэслэн хийгддэг бөгөөд хиймэл оюун, машин сургалттай хослон ажилдаг тэгсэнээр илүү нарийн үр дүнд хүч чаддаг. Computer vision чиглэлийн OpenCV Фреймворк нь 1999 оноос албан ёсоор нээлтээ хийж хөгжүүлэгдэж эхлэсэн хамгийн том фреймворк юм. eyeGrade програм нь хамгийн олон хэрэглэгчдэд хүрсэн desktop програм юм. Mobile төхөөрөмжийн хувьд zipgrade гэх application-ыг гадаадад ашигладаг мөн kahook гээд арай өөр төрлийн илүү өргөн хүрээг хамарсан application маш их ашиглагддаг.
\subsection{Технологийн судалгаа}
	\subsubsection{React Native}
	RN нь javascript дээр бичигдсэн ios болог android application хөгжүүлэхд зориулагдсан hybrid фреймворк юм. RN нь react javascript дээр тулгуурладаг. React - ийн component level pattern нь системийг өргөтгөхөд илүү хялбар уян хатан болгодог.
	\subsubsection{OpenCV}
	Opencv нь с болон с++  дээр бичигдсэн мөн бүх төрлийн cpu-дээр ажилладаг, бараг бүх түгээмэл програмчлалын хэл дээр байдаг. Cannyedge болон otsi-ийн илрүүлэлтийн алгоритмуудыг зэрэг маш олон хэрэгцээ ихтэй сангуудыг агуулдаг
	\subsubsection{MongoDB}
	MongoDB өгөгдөлийн сан нь харьцаа өгөдөлийн сан биш бөгөөд өгөдөлийн давхардалт зэрэг зүйлийг шалгадаггүй. Давуу тал нь илүү хурдан тооцоолох чадвартай. Бүх өгөгдөл collection-д хадгалагддаг.
	\subsubsection{Python}
	Python нь бүтээмж өндөртэй програмчлалын хэл юм.
	\subsection{Бүлгийн дүгнэлт}
	Python нь OpenCV-тэй сайн зохицож ажилладаг бөгөөд бичиглэл ойлгомжтой гэдэг үйднээс сонгосон. MongoDB бол харьцаа өгөгдөлийн сангаас хамаагүй уян хатан өөрчлөгдөх боломж хялбар бөгөөд өргөтгөхөд хялбар энгийн учир сонгосон. React-native ч мөн адил өргөтгөх өөрчлөх илүү хялбар гэдэг үүднээс сонгосон.

% Chapter Template

\section{Хөгжүүлэлтэд ашиглах технологи, хэрэгслүүд} % Main chapter title

\label{Chapter2} % Change X to a consecutive number; for referencing this chapter elsewhere, use \ref{ChapterX}

%----------------------------------------------------------------------------------------
%   SECTION 1
%----------------------------------------------------------------------------------------
\section{Технологийн тухай}
    \subsection{Технологи, прэймворк болон АПИ -ийн тухай үндсэн ойлголт}
    \subsection{Тэдгээрийн давуу ба дутагдалтай талууд}
    \subsection{Бусад ижил төстэй технологи, прэймворк, АПИ ба тэдгээртэй харьцуулсан судалгаа}
    \subsection{Технологийн сонголт дээрх SWOT шинжилгээ}
\section{Технологийн ерөнхий шаардлага}
    \subsection{Үйлдлийн системийн шаардлага}
    \subsection{Програмчлалын хэлний шаардлага}
    \subsection{Суулгах програмчлалын хэрэгслүүд}
    \subsection{Өгөгдлийн сангийн шаардлага}
\section{Суурилуулах заавар}
    \subsection{Сервер талд суурилуулах заавар}
    \subsection{Клиент талд суурилуулах заавар}
    \subsection{Програмчлалын хэрэгсэлтэй холбох заавар}
\section{Функц үйлдлүүд, тэдгээрийг хэрэглэх тухай}
%Класс, функц, массив, өгөгдлийн сантай холбох, өгөгдлийн сангаас өгөгдөл авах, бичих, устгах ... гэх мэтийг нэг бүрчлэн хэрхэн хэрэглэж байгааг тайлбарлах, жишээ гаргах
    \subsection{...}
    \subsection{...}
% Бүлэг 3

\section{\ttitle-н шинжилгээ} % Зарим нэг зөвлөмж

\label{Chapter3} % Энэ бүлэг рүү ишлэл хийх бол \ref{Chapter3} командыг ашигла 
%-------------------------------------------------------------------------------
%	SECTION
%-------------------------------------------------------------------------------

\section{Системийн үйл ажиллагааны тухай}

%	SECTION
%-------------------------------------------------------------------------------
\section{Системийг ашиглах хэрэглэгчид}
  \subsection{Системийн оролцогч}
  \subsection{Системийн тоглогч}
  
%	SECTION
%-------------------------------------------------------------------------------
\section{Функцийн шаардлага}

%	SECTION
%-------------------------------------------------------------------------------
\section{Функцийн бус шаардлага}

%	SECTION
%-------------------------------------------------------------------------------  
\section{Юзкейс диаграм}

%	SECTION
%-------------------------------------------------------------------------------  
\section{Юзкейс тодорхойлолт}

%	SECTION
%-------------------------------------------------------------------------------
\section{Шинжилгээний класс диаграм}

%	SECTION
%-------------------------------------------------------------------------------
\section{Шинжилгээний дарааллын диаграм}

%	SECTION
%-------------------------------------------------------------------------------
\section{Үйл ажиллагааны диаграм}

%	SECTION
%-------------------------------------------------------------------------------
\section{Бүлгийн дүгнэлт}

\pagestyle{plain}
% Ерөнхий дүгнэлт

\phantomsection
\addchaptertocentry{Ерөнхий дүгнэлт}
\label{Summary} % Энэ бүлэг рүү ишлэл хийх бол
%-------------------------------------------------------------------------------
%	Summary
%-------------------------------------------------------------------------------
\section*{Ерөнхий дүгнэлт}



\end{document}